\chapter{Implementácia}
\label{kap:implementation}
TODO

\section{Použité technológie}
Nová verzia pôvodnej aplikácie je plne implementovaná v jazyku TypeScript. Používateľské rozhranie je vytvorené pomocou knižnice React a na správu globálneho stavu sa používa knižnica Redux (všetky sú opísané v podkapitole \ref{sec:technologie}). Prirodzene sme sa teda rozhodli pokračovať v používaní týchto technológií aj pri tvorbe našich rozšírení. TypeScript sa ukázal ako veľmi nápomocný aj pri oboznamovaní sa s existujúcou implementáciou, keďže vďaka definovaným typom bolo výrazne jednoduchšie porozumieť očakávaným tvarom objektov alebo argumentov funkcií.

Tiež bolo potrebné vybrať knižnicu, pomocou ktorej budú implementované grafové editory. V kontexte už spomenutých technológií sme ako najvhodnejšie riešenie zvolili knižnicu ReactFlow. Na základe môjho hľadania išlo o jedinú knižnicu, ktorá podporovala už spomenuté technológie a zároveň poskytovala dostatočné množstvo konfigurácie na implementáciu všetkých navrhnutých funkcionalít.

\section{Refaktorizácia stavu pôvodnej aplikácie}

\section{Integrácia nových nástrojov}

\subsection{Organizácia}
Súborovú štruktúru pôvodnej aplikácie sme v zásade nemenili. Držali sme sa stanovenej konvencie, podľa ktorej je každá funkcionalita organizovaná v priečinku \texttt{src/features}. Zdrojové súbory pre jednotlivé nástroje preto ukladáme do samostatných podpriečinkov pomenovaných podľa daného nástroja. Obsah týchto priečinkov je rôzny, no každý obsahuje aspoň hlavný React komponent pre daný nástroj a súbor definujúci jeho Redux stav. V rámci stavu sú spravidla implementované reducery a potrebné selektory.

\subsection{Integrácia komponentov editorov}
Predošlá verzia aplikácie používala na zobrazovanie interpretácií textové vstupy, tie boli obalené v komponente \texttt{InterpretationInput} a následne používané naprieč aplikáciou. My máme však viac komponentov pomocou ktorých je možné interpretáciu konkrétneho symbolu reprezentovať.

Tento problém rieši nový komponent \texttt{InterpretationEditor}, ktorý zaobaľuje všetky editory vrátane pôvodného komponentu. Jedným z jeho vstupných parametrov je názov symbolu, ktorého interpretáciu reprezentuje. Na základe tohto názvu získa zo stavu typ editora, ktorý sa má zobraziť. Ak ide o textový editor zobrazuje sa pôvodný komponent \texttt{InterpretationInput}. V opačnom prípade sa vykreslí nový komponent \texttt{DrawerEditor}.

\texttt{DrawerEditor} integruje spoločné rozhranie editorov, ktoré sa skladá z ďalších troch častí. Jedným z nich je komponent panela nástrojov \texttt{EditorToolbar}. Ten implementuje rozhranie filtrovania. Ďalším je \texttt{EditorError}, ktorý predstavuje chybový banner. Treťou časťou je samotný komponent editora, ktorý je určený na základe získaného typu.

\subsection{Synchronizácia stavu}
Pri synchronizácie stavu jednotlivých editorov je potrebné, aby prebiehala automaticky. Teda každá zmena vykonaná v jednom editore sa musí okamžite premietnuť aj do stavu ostatných editorov. Každý z editorov však reprezentuje interpretáciu konkrétneho symbolu vo svojom stave iným spôsobom, čo predstavuje problém.

Aktuálna interpretácia konkrétneho symbolu je uložená v rámci stavu štruktúry a je možné ju meniť pomocou príslušných akcií:
\begin{itemize}
    \item \texttt{updateInterpretationPredicates} - akcia na aktualizáciu interpretácie predikátu
    \item \texttt{updateInterpretationFunctions} - akcia na aktualizáciu interpretácie funkcie
\end{itemize}
Súčasťou týchto akcií je okrem jej typu aj informácia o názve daného symbolu ako aj samotná aktualizovaná interpretácia.

Každý editor má v súvislosti s týmito akciami dve úlohy. Prvou je vyvolať príslušnú akciu v prípade, že sa nejakým spôsobom v jeho pohľade zmenila samotná interpretácia. To vyžaduje okrem iného aj konverziu medzi stavom reprezentácie a stavom ukladaným v štruktúre.

Druhou je vo svojom reducery reagovať aj na tieto akcie a na základe informácií v nich obsiahnutých aktualizovať svoj vnútorný stav. Na to bolo potrebné využiť špeciálnu funkcionalitu knižnice Redux, pomocou ktorej je možné v jednom reducery reagovať na akcie iného.

Ukážka takéhoto reducera (\ref{prog:extra_red}) demonštruje implementáciu tejto funkcionality pre Databázový editor pri zmene interpretácie predikátu. Každý reducer sa definuje v rámci funkcie \texttt{extraReducers}, ktorá dostáva ako argument objekt \texttt{builder}. Pomocou metódy \texttt{addCase} je možné špecifikovať akciu (prvý argument) pri ktorej sa má príslušný reducer zavolať (druhý argument).


\begin{lstlisting}[float,language=JavaScript,caption={Funkcia používaná na generovanie hierarchie},label=prog:extra_red]
extraReducers(builder) {
  builder.addCase(updateInterpretationPredicates, (state, action) => {
    if (action.meta.source === "databaseView") return;

    const { key, value } = action.payload;
    syncInterpretation(key, "predicate", value, state);
  });
}
\end{lstlisting}

\section{Implementácia grafových editorov}
TODO

%\subsection{Štruktúra}
%Implementácia grafových editorov je organizovaná v rámci priečinku \texttt{features/graphView}, ktorý je ďalej členený do logických celkov nasledovne:

%\begin{itemize}
%\item \texttt{/graphViewComponent} - obsahuje spoločný komponent pre zobrazovanie grafov, do ktorého sú jednotlivé reprezentácie vkladané
%\item \texttt{/graphs} - priečinok obsahuje hlavne zdrojové súbory jednotlivých reprezentácii organizovaných do podpriečinkov
%\begin{itemize}
%\item \texttt{/BipartiteGraph}
%\item \texttt{/HasseDiagram}
%\item \texttt{/OrientedGraph}
%\item \texttt{/common} - obsahuje súbory so spoločnou konfiguráciou editorov
%\item \texttt{/graphComponents} - obsahuje zdrojové kódy jednotlivých komponentov grafov, ako napríklad vrcholy alebo hrany
%\end{itemize}
%\item \texttt{/helpers} - obsahuje rôzne pomocné funkcie
%\end{itemize}

\subsection{Štruktúra stavu}
Stav grafov predstavuje obyčajný objekt, v ktorom každý záznam uchováva stav jednotlivých reprezentácií pre každý binárny predikát alebo unárnu funkciu. Kľúče sú kombináciou názvu a typu príslušného symbolu. Hodnota asociovaná s konkrétnym kľúčom je opäť objekt s dvoma atribútmi. Prvým je \texttt{tupleType}, ten predstavuje typ symbolu, a teda môže nadobúdať hodnoty \texttt{'predicate'} alebo \texttt{'function'}. Druhým je atribút \texttt{state}, ktorý je ďalej členený podľa jednotlivých reprezentácií na \texttt{state.oriented}, \texttt{state.bipartite} a \texttt{state.hasse}. Konkrétny stav je tvorený atribútmi \texttt{nodes} a \texttt{edges}, tie predstavujú zoznam vrcholov a hrán.

Tvar konkrétneho vrcholu alebo hrany je určený samotnou knižnicou a nie je odporúčané ho priamo rozširovať o vlastné atribúty. Pre tento účel je vyčlenený atribút \texttt{data}, ktorého tvar možno ľubovoľne špecifikovať. V prípade vrcholov ho využívame na zaznamenanie stavu, v ktorom sa vrchol nachádza pomocou atribútov \texttt{error}, \texttt{ghost}, \texttt{hatched} a \texttt{leftover}, všetky z nich sú voliteľné a typu \texttt{boolean}. Vrchol pre Bipartitný graf obsahuje ešte aj atribút \texttt{origin}, vyjadrujúci jeho príslušnosť do skupiny, môže nadobúdať hodnoty \texttt{'domain'} alebo \texttt{'range'}.

\subsection{Práca so stavom}
Jednotlivé dáta zo stavu sú komponentom grafov sprístupnené pomocou selektorov. To dovoľuje oddeliť transformáciu dát od implementácie používateľského rozhrania. Jednotlivé selektory je navyše možné kombinovať, čo uľahčuje ich použiteľnosť.

Napríklad úlohou selektora \texttt{selectNodes} je pre konkrétnu reprezentáciu získať zoznam vrcholov, aplikovať na ne aktuálne filtre domény a výsledok vrátiť v podobe zoznamu. Ako vstupné argumenty vyžaduje typ symbolu interpretácie, jeho názov a navyše ešte aj typ samotnej reprezentácie. Z týchto informácií je schopný zo stavu vybrať príslušný zoznam vrcholov. Na získanie domény po aplikovaní všetkých filtrov sa používa ďalší selektor \texttt{selectRelevantDomain}. Nakoniec je ešte potrebné zo zoznamu vrcholov odstrániť tie, ktoré nereprezentujú prvok nachádzajúci sa vo vyfiltrovanej doméne.

Na podobnom princípe funguje aj selektor \texttt{selectEdges}. Jeho úlohou je vrátiť zoznam hrán, ktorý sa dá rovno použiť pri zobrazení grafu. V prípade grafu reprezentujúceho hasseho diagramu to napríklad znamená potrebu odstrániť všetky nepotrebné hrany (podľa pravidiel opísaných v sekcií \ref{sec:reprezentacie-hasse}), keďže v jeho stave sú uložené hrany reprezentujúce celú interpretáciu.

Aktualizácia stavu prebieha pomocou reducerov a k nim prislúchajúcich akcií. Akcie na vyjadrenie zmien v stave vrcholov a hrán sú \texttt{nodesChanged} a \texttt{edgesChanged}. Zmena môže znamenať napríklad posunutie vrcholu alebo zakliknutie hrany. Pri vytvorení novej hrany je vyvolaná akcia \texttt{nodesConnected}, ktorej súčasťou sú aj identifikátory oboch vrcholov. Pomocou akcie \texttt{leftoverDeleted} je možné odstrániť chybný vrchol. Vyvoláva sa pri kliknutí na tlačidlo koša príslušného vrchola, popísaného v sekcií venovanej návrhu vrcholov \ref{sec:basic_graph}.

Dôležitou akciou je aj \texttt{syncGraphView}, ktorá slúži na inicializáciu stavu všetkých grafových reprezentácií všetkých binárnych predikátov a funkčných symbolov. Používa sa najmä pri synchronizácií grafových editorov so samotnou štruktúrou.

%Iba zacnem ze sa ukladaju hentak a blablabla. to uy je popisane v navrhu. potom spomeniem reducery a pluginy, mozno vypisat jednotlive funkcie pluginov a popisat ich

\subsection{Komponenty}
Hlavný komponent, ktorý obaluje komponenty konkrétnych reprezentácii je \texttt{GraphView}. Jeho úlohou je na základe poskytnutého grafového typu vykresliť správny komponent reprezentácie. Zobrazenie každej reprezentácie je implementované vrámci jej vlastného komponentu, konkrétne ide o \texttt{OrientedGraph}, \texttt{BipartiteGraph} a \texttt{HasseDiagram}. Tieto komponenty následne obaľujú komponent \texttt{ReactFlow} poskytovaný knižnicou, ktorý sa  využíva na vykreslenie samotného rozhrania grafu.

Úlohou jednotlivých komponentov je získavať potrebné dáta na vykreslenie grafu zo stavu príslušnej reprezentácie, ako sú napríklad vrcholy a hrany, prostredníctvom vyššie spomínaných selektorov. Zároveň sa v nich definuje konfigurácia komponentu \texttt{ReactFlow} špecifická pre daný typ grafu. Príkladom možnosti, ktorú knižnica pri konfigurácii ponúka je funkcia \texttt{isValidConnection}. Tú po jej špecifikovaní knižnica volá pri každom pokuse o vytvorenie novej hrany, čo poskytuje možnosť implementovania vlastnej validačnej logiky pre vytváranie hrán. Táto funkcionalita sa využíva napríklad v komponente \texttt{HasseDiagram}, kde sa kontroluje, či by pridaním hrany ostali neporušené všetky podmienky čiastočného usporiadania.

Menšie komponenty \texttt{PredicateNode} a \texttt{DirectEdge} predstavujú konkrétne vrcholy a hrany rozhrania grafu. O ich vykresľovanie sa stará samotná knižnica a taktiež je ich potrebné definovať vrámci konfigurácie. Pomocou komponentu \texttt{PredicateNode} je implementovaná napríklad funkcionalita vytvárania hrán alebo zobrazovania panelu s farbami unárnych predikátov.

%jednotlive komponenty ziskavaju data zo stavu
%spomenut nieco o konfiguraci atd.
%mozno pridat komponenty nodov edgov

%\subsection{Generovanie hierarchie}
%Na generovanie hierarchie pre Hasseho Diagrame sme sa rozhodli použiť knižnicu dagre.js \cite{dagrejs}. Jej hlavnými výhodami oproti ostatným alternatívam bola najmä jednoduchá konfigurácia a relatívne malá veľkosť balíčka.
%%
%Funkciu používanú na generovanie hierarchie možno vidieť na ukážke \ref{prog:hasse}. V prvom rade je potrebné vytvoriť inštanciu triedy \texttt{Graph}, ktorú poskytuje samotná knižnica. Následne sa pomocou metódy \texttt{setGraph} špecifikuje konfigurácia hierarchie. Nám z veľkej časti vyhovovala základná konfigurácie. Jediný atribút, ktorý bolo potrebné zmeniť bol \texttt{rankdir}. Ten určuje smer usporiadania vrcholov. V našom prípade chceme odspodu nahor. Po nastavení konfigurácie je potrebné špecifikovať jednotlivé vrcholy a hrany pomocou metód \texttt{setNode} a \texttt{setEdge}.
%%
%\begin{lstlisting}[float,language=JavaScript,caption={Funkcia používaná na generovanie hierarchie},label=prog:hasse]
%function computeLayoutHasse(nodes: Node[], edges: Edge[]) {
%  const dagreGraph = new dagre.graphlib.Graph();
%  dagreGraph.setGraph({ rankdir: "BT" });
%
%  nodes.forEach((n) => dagreGraph.setNode(n.id, {}));
%  edges.forEach((e) => dagreGraph.setEdge(e.source, e.target));
%
%  dagre.layout(dagreGraph);
%
%  return nodes;
%};
%\end{lstlisting}

%\section{Stav aplikácie}
%TODO
%\subsection{Štruktúra}
%TODO
%\subsection{Synchronizácia}
%TODO

\section{Integrácia s Logickým pracovným zošitom}
TODO
\subsection{Importovanie a exportovanie stavu}
TODO
%\subsection{Logický kontext}
%TODO

\section{Ďalšie vylepšenia Prieskumníka štruktúr}
TODO
\subsection{Refaktorizácia Henkinovej-Hintikkovej hry}
TODO